\documentclass[12pt,a4paper,oneside]{report}
\usepackage[utf8]{inputenc}
\usepackage{amsmath}
\usepackage{amsfonts}
\usepackage{amssymb}
\usepackage{booktabs}
\usepackage{graphicx}
\usepackage{hyperref}
\usepackage{geometry}
\geometry{margin=1in}

\title{\textbf{An Adaptive Framework for PCI DSS Implementation in Resource-Constrained Environments: The Case of Yemeni Banking Sector}}
\author{\textbf{Ahmed Allaw} \\ 
	Sana'a University, Information Systems Department}
\date{}

\begin{document}
	
	\maketitle
	
	\tableofcontents
	\newpage
	
	% ==========================================
	% CHAPTER ONE: INTRODUCTION
	% ==========================================
	\chapter{Introduction}
	
	\section{Introduction}
	
	\subsection{Digital Transformation in Banking}
	The contemporary global banking landscape has undergone a monumental and irreversible revolution, fundamentally altering its core operational structures through the widespread adoption of digital distribution channels. This seismic shift is characterized by the replacement of traditional, paper-based workflows with high-speed internet portals and versatile mobile banking applications. Globally, this transition is no longer a luxury but a fundamental survival mechanism for financial institutions aiming to remain competitive in a hyper-connected economy.
	
	\subsection{Growth of Electronic Banking in Yemen}
	Within the specific and complex socio-economic context of Yemen, the aggressive shift toward digitalization has been fueled by a pressing, multifaceted demand for banking services that prioritize streamlined efficiency, reduced overhead costs, and seamless accessibility. Historically, traditional physical banking institutions in Yemen have struggled to provide such agility due to bureaucratic hurdles, geographical constraints, and an aging infrastructure. Consequently, the promise of remote accessibility has gained immense traction among the Yemeni populace and businesses alike. Despite this demand, the adoption of electronic banking services in Yemen still lags significantly behind global benchmarks, primarily due to a unique blend of psychological barriers and systemic limitations.
	
	\subsection{Cybersecurity Challenges Associated with Digital Banking}
	However, this rapid transition into the digital realm is not without significant peril. The migration of financial assets and sensitive personal data to online servers has opened a multitude of technological vulnerabilities. The inherent susceptibility of the Yemeni financial infrastructure to sophisticated cyber-attacks—ranging from ransomware and phishing to advanced persistent threats—introduces a perilous environment. 
	
	In this digital frontier, the privacy and fundamental integrity of sensitive user data are under constant threat from both local and international malicious actors. Such security breaches often result in catastrophic financial losses, the potential collapse of smaller financial entities, and a total, irreversible erosion of institutional credibility. In a country already grappling with economic instability, the loss of public trust in the digital banking system could lead to a widespread return to informal, unregulated financial shadow systems, further crippling the formal economy.
	
	\section{Research Background}
	
	\subsection{Overview of PCI DSS}
	In direct response to escalating global threats, the Payment Card Industry Data Security Standard (PCI DSS) has established itself as the premier international framework dedicated to the rigorous safeguarding of cardholder data and the securing of the broader transaction lifecycle. Developed and maintained by the PCI Security Standards Council—a consortium of industry entities such as Visa, MasterCard, and American Express—this framework provides a universal language for protecting sensitive transaction data. The standard provides a comprehensive roadmap for technical and operational requirements designed to protect account data.
	
	\subsection{Importance of Payment Card Security}
	At its core, the goal of PCI DSS is to ensure cardholder data security through rigorous, worldwide standards that prevent fraud and data breaches. A fundamental concept within the standard is the Cardholder Data Environment (CDE), dictating that the standard applies to any entity that handles, stores, or transmits cardholder account data. To protect this data, the standard mandates key security controls, including firewall management to establish a robust perimeter, data encryption to ensure intercepted data remains unreadable, and strict access controls to restrict physical and digital access solely to authorized personnel.
	
	\subsection{Global Relevance of PCI DSS}
	For financial institutions worldwide, compliance is essential for establishing a baseline of security and trust. For a developing economy like Yemen, the adoption of PCI DSS is far more than a simple technical upgrade or a routine compliance exercise; it serves as a foundational strategic pillar and a badge of international legitimacy. It is the primary vehicle for winning the hard-earned confidence of a skeptical marketplace that remains acutely aware of high-level digital risks. Without the assurance provided by such rigorous standards, the growth of e-commerce and digital payments in Yemen will likely remain stunted.
	
	\section{Problem Statement}
	
	\subsection{Security Challenges Facing Yemeni Banks}
	While the advantages of digital banking are clear, the digital transformation of the banking sector remains a complex challenge in developing nations. In Yemen, the integration of digital finance faces a foundational barrier in the form of heightened risk perception. Users often associate digital platforms with technical vulnerabilities or the threat of cybercrime, leading to a preference for traditional, cash-based transactions. Furthermore, a ``Security Paradox'' exists: while privacy concerns are not the sole statistical deterrent to adoption, a deeper cultural reliance on established financial behaviors dictates that trust must be systematically built through visible and robust security measures.
	
	\subsection{Infrastructure Limitations}
	Implementing globally recognized standards like PCI DSS requires a stable operational environment, which Yemen currently lacks. The infrastructure suffers from severe deficits, specifically unreliable electricity and poor internet connectivity, which pose significant challenges to maintaining continuous security operations and real-time monitoring. Standard PCI DSS requirements are often designed for high-bandwidth, stable environments, making direct application in Yemen highly problematic.
	
	\subsection{Compliance Difficulties in Resource-Constrained Environments}
	Beyond physical infrastructure, the full implementation of PCI DSS is hindered by systemic and organizational roadblocks. There is a critical shortage of specialized cybersecurity talent and a lack of robust regulatory frameworks to enforce compliance. Additionally, the high financial cost of the standard audit and certification process presents a significant barrier for local banks operating under tight capital constraints. 
	
	Organizational readiness is also fundamentally lacking, with security processes being the least developed competency within Yemeni banks. Consequently, attempting to enforce the standard mandate of simultaneous compliance across all twelve PCI DSS requirements is entirely impractical for these institutions.
	
	\sectionsection{Research Aim}
	The primary aim of this research is to conduct an exhaustive, multidimensional evaluation of the existing security performance levels across Yemeni banking institutions and to propose a specialized, adaptive framework for the implementation of the Payment Card Industry Data Security Standard (PCI DSS) that accommodates the unique resource constraints and infrastructural limitations of the region.
	
	\sectionsection{Research Objectives}
	To fulfill the overarching aim, this study pursues the following measurable objectives:
	\begin{itemize}
		\item To evaluate the existing information security management systems and the behavioral factors affecting the acceptability of electronic banking in Yemen.
		\item To diagnose the multifaceted technical, systemic, and infrastructural challenges that complicate the full implementation of PCI DSS in the Yemeni context.
		\item To design an adaptive, phase-based compliance roadmap that prioritizes critical vulnerabilities—such as core data protection and infrastructure resilience—before pursuing full international certification.
		\item To formulate localized security strategies, incorporating AI-driven fraud detection for resource optimization and hybrid security awareness programs tailored to local cultural behaviors.
	\end{itemize}
	
	\sectionsection{Research Questions}
	Guided by the research objectives, this study seeks to answer the following questions:
	\begin{itemize}
		\item What is the current state of information security management systems and public acceptance regarding electronic banking within the Yemeni banking sector?
		\item What are the primary infrastructural, regulatory, and organizational barriers preventing the immediate implementation of standard PCI DSS requirements in Yemen?
		\item How can the stringent requirements of PCI DSS be restructured into a tiered, adaptive framework suitable for a low-bandwidth, resource-constrained environment?
		\item In what ways can artificial intelligence and localized human-centric processes compensate for the severe shortage of specialized cybersecurity personnel in Yemeni financial institutions?
	\end{itemize}
	
	\section{Significance of the Study}
	
	\subsection{Academic Significance}
	Academically, this study addresses critical research gaps regarding cybersecurity in developing nations. It contributes to the existing body of knowledge by proposing localized security maturity models specifically designed for the resource-constrained environments of Yemeni banks and Small and Medium Enterprises (SMEs). Furthermore, it explores the novel intersection of artificial intelligence and automated fraud detection systems operating within the boundaries of PCI DSS requirements, providing a theoretical foundation for maintaining rigorous oversight with limited human cybersecurity resources.
	
	\subsection{Practical Significance for Yemeni Banks}
	From a practical standpoint, overcoming the formidable barriers to standard compliance presents a transformative opportunity for the modernization of Yemen's financial architecture. The proposed adaptive framework allows banks to move beyond a mere ``check-the-box'' compliance exercise, instead building a resilient foundation. Successfully adopting these protocols aligns local institutions with global financial best practices, making them more attractive to international partners and foreign investment. 
	
	Ultimately, this framework will act as a catalyst for national financial inclusion, enabling the secure rollout of innovative services like micro-finance and secure international remittances, utilizing cybersecurity as an enabler of economic progress and national stability.
	
	\sectionsection{Scope of the Study}
	This research is strictly confined to the Yemeni banking sector and its transition toward electronic banking. The operational scope focuses explicitly on the implementation parameters of the Payment Card Industry Data Security Standard (PCI DSS). The study limits its systemic analysis to the immediate resource constraints faced within the country, specifically infrastructural deficits (power and internet limitations), organizational readiness, specialized staffing shortages, and localized consumer trust factors.
	
	\sectionsection{Structure of the Thesis}
	The thesis is organized to systematically navigate the complexities of securing digital finance in a developing economy. Following this introductory chapter, subsequent chapters will synthesize empirical research concerning the Yemeni banking environment and consumer behavior. The research will deeply analyze the technical architecture of PCI DSS, examining network segmentation, access controls, and data encryption requirements. It will then present the detailed diagnosis of infrastructural and regulatory limitations before transitioning into the comprehensive presentation of the proposed Adaptive PCI DSS Implementation Framework. The final sections will outline the tiered progression roadmap—from core data protection to full certification—and conclude with future directions for integrating AI automation and hybrid security awareness models into the sector's long-term strategy.
	
	% ==========================================
	% CHAPTER TWO: LITERATURE REVIEW
	% ==========================================
	\chapter{Literature Review}
	
	\section{Introduction to the Literature Review}
	The rapid evolution of information technology has fundamentally shifted the global financial paradigm, driving a comprehensive transition from traditional, brick-and-mortar banking operations to digitalized electronic banking, commonly referred to as e-banking. While this digital migration optimizes operational efficiency, it simultaneously introduces substantial vulnerabilities, exposing critical financial and organizational assets to diverse categories of cyber-attacks that directly threaten the confidentiality and integrity of sensitive transaction data. 
	
	In response to these escalating threats, financial institutions globally rely on standardized compliance frameworks to mitigate risk and preserve data integrity. This literature review provides an in-depth, systematic evaluation of the implementation of the Payment Card Industry Data Security Standard (PCI DSS), focusing on the operational realities of resource-constrained environments, specifically the Yemeni banking sector. 
	
	By synthesizing empirical research concerning existing information security management systems (ISMS) and evaluating the socio-technological factors that dictate consumer acceptance of e-banking, this chapter outlines the systemic challenges, regulatory impediments, and infrastructural limitations that characterize the Yemeni financial ecosystem. Furthermore, it explores the theoretical underpinnings of technology adoption under conditions of economic instability, highlighting how standardized security models function as a primary mechanism for building institutional trust and driving national financial modernization.
	
	\sectionsection{Electronic Banking in Developing Countries}
	
	\subsection{Evolution of e-banking}
	The modern banking landscape has experienced an irreversible structural revolution, changing its foundational operational paradigms through the deployment of digital distribution channels. This transition is defined by the systematic replacement of legacy, paper-based workflows with high-speed internet portals and multi-functional mobile banking applications. On a global scale, the adoption of these digital ecosystems has ceased to be an optional technological upgrade; rather, it has become a fundamental survival mechanism for financial entities seeking to maintain competitiveness within a hyper-connected economy.
	
	\subsection{Challenges in developing economies}
	Despite global advancements, the execution of digital transformation initiatives within the banking sectors of developing nations remains an exceptionally complex and uneven endeavor. The propagation and maturity of e-banking infrastructures within these regions frequently lag behind international benchmarks. This developmental delay is primarily driven by a multi-faceted combination of persistent psychological barriers among the consumer base and deep-seated systemic limitations within the domestic operating environments. Financial institutions in these contexts must manage technological integration while operating under severe resource constraints, unstable core infrastructure, and volatile market conditions.
	
	\subsection{Situation in Yemen}
	Within the specific socio-economic context of Yemen, the transition toward digital banking channels has been accelerated by a critical, localized demand for financial services capable of delivering streamlined efficiency, minimal overhead expenditures, and widespread accessibility. Historically, conventional physical banking entities across Yemen have failed to offer adequate operational agility. This institutional inertia stems from deep bureaucratic hurdles, severe geographical constraints, and an aging physical infrastructure. Consequently, the strategic promise of decentralized digital payments has achieved significant traction among both commercial enterprises and the general populace.
	
	However, this rapid digital migration introduces substantial operational peril. The transfer of financial assets and highly sensitive personally identifiable information (PII) to online servers has exposed a wide array of technological vulnerabilities. The Yemeni financial infrastructure exhibits a pronounced susceptibility to sophisticated cyber-attacks, ranging from ransomware and phishing campaigns to advanced persistent threats (APTs). This environment leaves the fundamental integrity and privacy of user data under constant threat from both domestic and international malicious actors. 
	
	Security failures within this context extend far beyond temporary technical disruptions; they routinely culminate in catastrophic financial losses, threaten the structural viability of smaller financial institutions, and cause an irreversible erosion of public credibility. In an environment already burdened by broader economic instability, the collapse of public confidence in digital banking risks triggering a mass migration back to informal, unregulated financial ``shadow'' systems, further destabilizing the formal national economy.
	
	\sectionsection{Cybersecurity in the Banking Sector}
	
	\subsection{Major cyber threats}
	The contemporary digital frontier exposes financial institutions to a hostile threat landscape where sensitive assets are targeted by a diverse set of threat vectors. Financial systems operating within developing or unstable jurisdictions face highly aggressive cyber-enabled operations. The primary technical threats identified within the banking sector encompass:
	\begin{itemize}
		\item \textbf{Ransomware:} Malicious cryptographic attacks designed to extort financial institutions by rendering core operational databases inaccessible.
		\item \textbf{Phishing:} Social engineering schemes executed by local and international actors to compromise privileged credentials and access sensitive financial networks.
		\item \textbf{Advanced Persistent Threats (APTs):} Highly sophisticated, prolonged cyber-campaigns directed by skilled threat actors aiming to infiltrate the core transactional environment and exfiltrate data over extended durations.
	\end{itemize}
	
	\subsection{Importance of information security management}
	Robust information security management serves as the primary mechanism preventing systemic operational collapse within modern financial systems. Weaknesses in security protocols do not merely present minor technical inconveniences; they directly imperil the financial solvency of banking firms and dissolve customer trust. Consequently, the systematic implementation of a structured Information Security Management System (ISMS) is vital. A mature ISMS provides the necessary governance framework to proactively identify vulnerabilities, enforce administrative accountabilities, and establish technical parameters that preserve the confidentiality, availability, and absolute integrity of transaction data.
	
	\subsection{Security requirements for financial institutions}
	To safeguard data assets effectively, financial institutions must move beyond reactive security measures toward proactive compliance regimes. These institutions require comprehensive technical and operational frameworks capable of securing the entire transaction lifecycle. Security models such as the PCI DSS are essential to this objective, acting as the foundational mechanism for generating trust within the consumer marketplace. Adherence to these strict standards ensures that data security is embedded directly into the technical architecture of the network, protecting the institution against both external breaches and internal process failures.
	
	\sectionsection{PCI DSS Overview}
	
	\subsection{History and purpose}
	The Payment Card Industry Data Security Standard (PCI DSS) represents the premier global framework dedicated to the rigorous defense of cardholder account data and the broader transactional environment. The standard was originally conceptualized, developed, and continues to be maintained by the PCI Security Standards Council (PCI SSC), an international consortium comprising major payment card providers including Visa, MasterCard, and American Express. 
	
	The overarching objective of this framework is to establish a universal, standardized language for technical and operational requirements designed specifically to mitigate payment card fraud, minimize data breaches, and secure sensitive financial transactions on a global scale.
	
	\subsection{Cardholder Data Environment (CDE)}
	A fundamental architectural concept governing the implementation of PCI DSS is the Cardholder Data Environment (CDE). The standard mandates that its technical and operational stipulations apply directly to any system component, network segment, or entity that stores, processes, or transmits cardholder account data or sensitive authentication information. The scoping criteria enforced by the standard are absolute: if an operational asset, device, or software application interacts with or touches a credit card number, it is automatically classified as ``in scope'' and must comply with the full spectrum of PCI DSS controls.
	
	\subsection{Network segmentation}
	Securing every individual endpoint and device across a sprawling, complex corporate enterprise network is frequently technically unfeasible and cost-prohibitive. To resolve this operational challenge, organizations deploy network segmentation as a primary scope reduction strategy. This practice involves the logical or physical partitioning of the corporate network into highly isolated and distinct security zones:
	
	\begin{table}[h!]
		\centering
		\caption{Network Zones Operational Definition and Security Status}
		\label{tab:network_zones}
		\begin{tabular}{p{5cm} p{10cm}}
			\toprule
			\textbf{Network Zone} & \textbf{Operational Definition \& Security Status} \\
			\midrule
			The Cardholder Data Environment (CDE) & The highly secure ``inner sanctum'' of the corporate architecture where all sensitive cardholder data resides, is processed, or is transmitted. \\
			\midrule
			Shared Network & Network zones containing systems that directly interact or communicate with the CDE, but are prohibited from storing sensitive account data. \\
			\midrule
			Corporate LAN & General corporate office systems (such as Human Resources, marketing, or general administrative departments) that must remain strictly isolated from the CDE. \\
			\bottomrule
		\end{tabular}
	\end{table}
	
	\subsection{Security controls}
	The structural integrity of the PCI DSS framework is maintained by several critical pillars of technical defense that must be continuously enforced:
	\begin{itemize}
		\item \textbf{Firewall Management:} The implementation of robust perimeter defenses designed to monitor, filter, and strictly control incoming and outgoing network traffic based on predefined security policies.
		\item \textbf{Data Encryption (Requirement 4):} The mandatory enforcement of cryptographic protocols to ensure that even if account data is intercepted during transmission across open, untrusted, or public networks, it remains unreadable and useless to unauthorized actors.
		\item \textbf{Access Control (Requirement 9):} The restriction of both physical and digital access to the cardholder data environment, ensuring that only explicitly validated personnel with a legitimate, documented business ``need to know'' can interact with secure systems.
	\end{itemize}
	
	\subsection{The twelve PCI DSS requirements}
	To attain full compliance validation, financial entities must successfully address six overarching security goals through the deployment of twelve highly specific technical and operational requirements:
	\begin{itemize}
		\item \textbf{Goal 1: Build and Maintain Secure Networks}
		\begin{itemize}
			\item Requirement 1 \& 2: Install and maintain robust firewall configurations to protect cardholder data, and systematically change all vendor-supplied default passwords and security parameters prior to system deployment.
		\end{itemize}
		\item \textbf{Goal 2: Protect Cardholder Data}
		\begin{itemize}
			\item Requirement 3 \& 4: Protect stored cardholder data through cryptographic storage mechanisms, and strictly encrypt all data during transmission across open, public networks.
		\end{itemize}
		\item \textbf{Goal 3: Maintain a Vulnerability Management Program}
		\begin{itemize}
			\item Requirement 5 \& 6: Deploy and regularly update anti-virus software, and develop and maintain highly secure systems and software applications.
		\end{itemize}
		\item \textbf{Goal 4: Implement Strong Access Control Measures}
		\begin{itemize}
			\item Requirement 7, 8 \& 9: Restrict digital access to cardholder data strictly by business ``need to know,'' assign unique identifiers to every individual with computer access, and strictly track and restrict physical access to cardholder data assets.
		\end{itemize}
		\item \textbf{Goal 5: Regularly Monitor and Test Networks}
		\begin{itemize}
			\item Requirement 10 \& 11: Track and monitor all access to network resources and sensitive cardholder data, and regularly test security systems and operational processes.
		\end{itemize}
		\item \textbf{Goal 6: Maintain an Information Security Policy}
		\begin{itemize}
			\item Requirement 12: Maintain a comprehensive policy that explicitly addresses information security responsibilities for all personnel across the organization.
		\end{itemize}
	\end{itemize}
	
	\section{Factors Affecting E-Banking Adoption in Yemen}
	
	\subsection{Risk perception}
	A primary psychological barrier impeding the widespread integration and acceptance of digital financial services among consumers in Yemen is a heightened perception of risk. Potential users frequently exhibit a widespread fear of experiencing direct financial loss through online channels. Consumers heavily associate digital banking interfaces with critical technical vulnerabilities and the persistent threat of targeted cybercrime. This elevated risk perception deters users from adopting digital methods, reinforcing a strong behavioral preference for traditional, physical, cash-based transactions.
	
	\subsection{Trust}
	The establishment of institutional trust is a prerequisite for achieving technology acceptance within a highly skeptical marketplace. Consumers within the Yemeni market remain acutely conscious of high-level digital risks, operational vulnerabilities, and the widespread potential for financial fraud. In the absence of formal, explicit assurances provided by verified, internationally recognized security models like PCI DSS, the broader expansion of e-commerce frameworks and digital payment ecosystems will remain severely stunted. Trust acts as the primary psychological bridge required to convert a risk-averse population into active participants in digital finance.
	
	\subsection{Awareness}
	The sluggish adoption rate of digital financial platforms is further compounded by substantial trust and awareness gaps across the population. Empirical data reveals a notable deficiency in public knowledge regarding the explicit, practical advantages of utilizing e-banking systems. Without a clear and comprehensive public understanding of how these technological tools simplify daily financial management and improve security, potential consumers remain unmotivated to migrate away from physical banking branches.
	
	\subsection{Usability}
	System usability represents a powerful determinant governing the behavioral intention to adopt e-banking applications. When digital platforms present intuitive interfaces and low operational complexity, user willingness to engage with the technology increases. However, within the Yemeni ecosystem, the positive impacts of usability are constrained by systemic structural challenges, meaning that user-friendly design alone is insufficient to drive complete national adoption in the absence of stable infrastructure and security assurances.
	
	\subsection{Consumer behavior and the ``Security Paradox''}
	The empirical evaluation of consumer behavior within Yemen reveals a highly unexpected phenomenon known as the Security Paradox. While conventional information security theories dictate that data privacy and technical security concerns serve as the primary drivers of user rejection or acceptance, localized data from Yemen contradicts this assumption. 
	
	Empirical analysis demonstrates that while interface usability and general public awareness are highly significant drivers of technology adoption, ``security and privacy'' metrics surprisingly do not exert a statistically significant impact on a consumer's ultimate behavioral decision to engage with e-banking platforms.
	
	This paradox indicates that user resistance is not simply a rational response to technical safety deficits. Instead, it reflects a deep cultural reliance on long-established financial behaviors and traditional social structures. Yemeni consumers assume a baseline level of insecurity due to the surrounding environment, shifting the adoption decision to factors of immediate utility and usability. 
	
	Therefore, implementing robust standards like PCI DSS is not necessarily required to satisfy immediate consumer demands for technical privacy metrics, but is essential for creating the structural stability and institutional legitimacy needed to gradually reshape deep-seated cultural financial behaviors.
	
	\section{PCI DSS Implementation Challenges in Yemen}
	
	\subsection{Technical limitations}
	The practical engineering tasks required to build a compliant banking architecture face extensive technical limitations within the local environment. Implementing advanced security components—such as Virtual Local Area Networks (VLANs) and Access Control Lists (ACLs) to partition the CDE—requires continuous administrative configuration. Furthermore, enforcing strong remote identity management through Authentication, Authorization, and Accounting (AAA) systems and Terminal Access Controller Access-Control System Plus (TACACS+) protocols requires highly stable network handshakes that are technically difficult to sustain in resource-constrained environments.
	
	\subsection{Infrastructure constraints}
	The broader operational environment in Yemen is characterized by severe, systemic infrastructure deficits that present persistent obstacles to maintaining basic information security operations. The nation’s financial institutions grapple with frequent power cuts, an unreliable electrical grid, and a lack of proper, consistent access to baseline computing technologies. Furthermore, highly volatile internet connectivity and deliberate internet throttles frequently disrupt real-time network tracking. Because standard compliance regimes assume high-bandwidth, uninterrupted, real-time telemetry data, these infrastructure constraints make standard, continuous compliance monitoring impossible without localized modifications.
	
	\subsection{Regulatory challenges}
	The regulatory landscape in Yemen introduces additional layers of complexity. Current national banking policies and central legal frameworks are outdated and fail to accommodate critical environmental, technological, and cultural realities. There is a pronounced absence of robust regulatory requirements and state-enforced legal mandates designed to compel financial institutions to systematically pursue and maintain compliance. Without a centralized, proactive regulatory authority to enforce clear compliance baselines, local banks lack the legal pressure required to prioritize costly security compliance programs over short-term operational survival.
	
	\subsection{Organizational readiness}
	Yemeni banking institutions exhibit low levels of organizational readiness, which hinders the successful adoption of advanced cybersecurity frameworks. Empirical assessments demonstrate that the ``process'' dimension exerts the most significant statistical influence over the improvement of security management within Yemen. However, this exact dimension currently exists in the least developed form within local financial entities. Organizational readiness is severely constrained by several factors:
	\begin{itemize}
		\item \textbf{Capital Constraints and High Costs:} The substantial financial expenditures required to complete formal external security audits and achieve official compliance certification present an immense barrier for local banks operating under highly restricted capital constraints.
		\item \textbf{Human Capital Shortages:} There is an acute, chronic shortage of specialized cybersecurity talent and experienced incident response professionals within the domestic labor market.
		\item \textbf{Policy Gaps:} The majority of active banks have completely failed to conceptualize or incorporate localized organizational policies capable of bridging the gap between technical requirements and actual human behavior.
	\end{itemize}
	
	\subsection{Security awareness issues}
	The organizational vulnerability of the financial sector is further heightened by pervasive security awareness deficiencies across multiple strata of the banking ecosystem. There is an immediate, critical requirement to elevate the baseline security literacy of both internal bank staff and external customers. In the absence of comprehensive, ongoing educational programs, users remain highly susceptible to sophisticated social engineering campaigns, credential harvesting, and phishing tactics that can bypass advanced technical perimeter controls.
	
	\section{Benefits of PCI DSS Adoption}
	
	\subsection{Trust enhancement}
	Achieving formal validation against the PCI DSS framework provides a powerful ``seal of security'' that fundamentally transforms the market standing of a Yemeni bank. This international certification functions as a visible strategic credential that effectively mitigates heightened consumer risk perceptions. By demonstrating a verified adherence to international security mandates, banks can successfully cultivate public confidence, overcome consumer skepticism, and actively accelerate the long-term adoption of electronic banking services.
	
	\subsection{Security improvement}
	From an engineering perspective, pursuing compliance allows an institution to significantly strengthen and mature its broader Information Security Management System (ISMS). By methodically executing all twelve requirements, a bank elevates its technical controls to match international best practices, aligning its operations with reputable frameworks such as ISO 27001. The practical benefits of this alignment include:
	\begin{itemize}
		\item \textbf{Robust Access Auditing:} Forcing the comprehensive implementation of system monitoring tools to record, track, and log all user access, ensuring rapid detection and response to unauthorized access attempts.
		\item \textbf{Network Isolation:} Enhancing data protection through the systematic deployment of VLANs and ACLs to physically and logically isolate cardholder data from general business traffic.
		\item \textbf{Cryptographic Communication:} Mandating that all internal and external servers handling sensitive authentication info communicate exclusively via secure cryptographic codes, ensuring data remains unreadable if intercepted.
	\end{itemize}
	
	\subsection{Support for digital transformation}
	The structural adoption of PCI DSS serves as a transformative enabler and catalyst for comprehensive national financial modernization. Operating within a fully validated, compliant environment allows banks to leverage cutting-edge financial technologies, such as Internet of Things (IoT) payment systems, which cannot be deployed securely in uncompliant settings. 
	
	Furthermore, it provides the secure foundation necessary to safely roll out innovative, high-impact services including mobile wallets, localized micro-finance initiatives, and secure international remittance channels. Ultimately, this security baseline grants international legitimacy, aligning local systems with global financial networks and making Yemeni institutions highly attractive to foreign investment and international banking partners.
	
	\section{Research Gap}
	
	\subsection{Gaps identified in the paper}
	The foundational body of global literature regarding cybersecurity compliance exhibits a pronounced geographical and situational bias. Standard compliance models are built upon the implicit assumption that the implementing organization operates within an optimized environment characterized by stable physical utilities, high-bandwidth connectivity, unlimited capital reserves, and a deep pool of specialized cybersecurity professionals. The existing literature fails to address how resource-constrained institutions can maintain compliance when these foundational assumptions are entirely absent.
	
	\subsection{Need for localized frameworks}
	A critical research gap exists regarding the formulation of localized security maturity models specifically engineered to accommodate the resource-constrained realities of small-to-medium enterprises (SMEs) and banking institutions operating within conflict-affected or developing regions. The standard mandate requiring immediate, simultaneous enforcement of all twelve requirements is structurally unviable in these environments. 
	
	Academic and practical research must therefore focus on designing tiered, phase-based progression roadmaps that prioritize critical data protection objectives first, allowing institutions to balance high audit costs against severe infrastructure deficits.
	
	\subsection{Role of AI and automation}
	There is a stark deficiency in empirical studies exploring the integration of Artificial Intelligence (AI) and automated compliance systems within human-capital-starved cybersecurity sectors. While standard frameworks rely on large, dedicated security operations teams to perform continuous log analysis, future research must address how AI-driven fraud detection systems and automated compliance monitoring tools can be deployed to execute real-time threat detection. Utilizing automation to fulfill monitoring mandates (such as Requirement 10) can effectively offset severe domestic shortages of specialized human expertise.
	
	\subsection{Security awareness integration}
	Finally, current models lack comprehensive integration mechanisms that simultaneously combine human factors, organizational workflows, and technical parameters. There is a distinct need to develop hybrid security awareness frameworks that incorporate a unified ``people, process, and technology'' approach tailored to unique local banking behaviors. 
	
	Research must focus on creating hybrid training modules that merge technical PCI DSS requirements with anti-social engineering strategies, alongside localized policy templates that translate abstract international regulations into practical administrative workflows that fit the existing bureaucratic realities of developing nations.
	
	\subsection{Chapter Summary}
	In summary, the academic literature establishes that the structural transition from legacy, paper-based workflows to digital electronic banking channels is an unavoidable prerequisite for economic survival and modernization. However, within developing and resource-constrained environments like Yemen, this rapid technological migration introduces severe cyber-risks, exposing highly sensitive financial assets to destructive threat vectors including ransomware, phishing, and advanced persistent threats. 
	
	The societal response to this transition is governed by a complex Security Paradox, wherein consumer technology acceptance is driven heavily by interface usability and general awareness, while technical security metrics fail to show a statistically significant impact due to a deep cultural reliance on traditional, cash-based financial behaviors. 
	
	To overcome this skepticism and prevent a total retreat into unregulated financial shadow systems, the adoption of strict security frameworks like the PCI DSS is imperative. Yet, the full execution of this standard is blocked by immense technical limitations, severe infrastructure deficits, a lack of cohesive regulatory enforcement, and low organizational readiness characterized by severe capital constraints and a shortage of specialized talent. 
	
	Because standard compliance paradigms assume a highly continuous, high-bandwidth environment managed by specialized security teams, a significant research gap remains regarding how to successfully execute these controls in volatile contexts. Resolving this gap requires moving away from rigid, uniform mandates toward localized, adaptive frameworks that utilize phased implementation tiers, automated AI monitoring tools, and culturally integrated training processes to build a secure, resilient financial foundation capable of driving national economic progress.
	
	% ==========================================
	% CHAPTER THREE: METHODOLOGY AND DEVELOPMENT
	% ==========================================
	\chapter{Research Methodology and Framework Development}
	
	\section{Introduction}
	The transition of the Yemeni financial sector from conventional, paper-based workflows to digital and electronic banking platforms represents a critical evolution driven by the demands of a hyper-connected global economy. While this digital shift optimizes organizational efficiency, minimizes operational overhead, and enhances user accessibility, it exposes the underlying financial infrastructure to complex cybersecurity vulnerabilities and sophisticated threat vectors. Securing sensitive transaction environments within a developing country like Yemen requires structured, systemic compliance mechanisms such as the Payment Card Industry Data Security Standard (PCI DSS). 
	
	However, direct implementation of standardized global security protocols is routinely obstructed by localized socio-economic, technical, and operational resource constraints. This chapter provides a comprehensive examination of the research methodology employed to study these dynamics, details the structural realities of the Yemeni e-banking environment, outlines the baseline technical architecture required for transactional data security, and develops a specialized, adaptive framework engineered specifically to facilitate phased PCI DSS compliance within resource-constrained financial institutions.
	
	\sectionsection{Research Methodology}
	The operational foundation of this study is built upon an inductive, multi-dimensional analysis tailored to examine the structural and behavioral parameters governing information security within the Yemeni banking sector. Given the structural challenges present within the target environment, the methodology avoids purely theoretical modeling in favor of a specialized approach grounded in the empirical operational realities of local financial institutions.
	
	\subsection{Analytical and Synthesis-Based Approach}
	The research methodology is systematically executed through an in-depth synthesis of recent empirical research addressing the technical architecture, regulatory limitations, organizational structures, and consumer behavioral metrics characteristic of the Yemeni financial ecosystem. By cross-examining distributed empirical data points, the analytical process targets the intersection of technology acceptance models, information security management frameworks, and the unique socio-economic disruptions local to Yemen. 
	
	This synthesis methodology isolates conflicting variables—such as technical safety versus consumer risk perception—to establish a data-driven baseline for architectural modifications. The research translates qualitative and quantitative findings from existing operational studies into a formalized, technically actionable security framework.
	
	\subsection{Framework Evaluation Parameters}
	The evaluation of the PCI DSS implementation landscape within the Yemeni banking sector is structured around four primary research parameters:
	\begin{itemize}
		\item \textbf{Infrastructure Assessment:} Auditing the current state of critical baseline dependencies, specifically power availability, network connectivity, and transmission stability, to establish realistic operational boundaries for security controls.
		\item \textbf{Regulatory and Policy Auditing:} Analyzing the presence, absence, and structural design of organizational security policies and national legal frameworks necessary to legally mandate and operationally enforce compliance metrics.
		\item \textbf{Organizational Readiness Modeling:} Evaluating internal process maturities, resource allocations, and administrative workflows inside local banking institutions to identify core structural vulnerabilities in security governance.
		\item \textbf{User Behavior Analysis:} Examining the psychological pillars, awareness deficits, and cultural dependencies of both banking personnel and consumers to map the systemic barriers impeding digital adoption and driving risk perception.
	\end{itemize}
	
	\section{Analysis of the Yemeni Banking Environment}
	To formulate a functional implementation roadmap, the specific operational realities, technical deficiencies, and behavioral paradoxes defining the contemporary Yemeni financial ecosystem must be thoroughly contextualized.
	
	\subsection{Current Cybersecurity Landscape}
	The aggressive migration of financial assets and critical cardholder profiles from physical registries to online servers has created an environment highly vulnerable to advanced cyber-attacks. The Yemeni digital frontier operates under continuous threat from local and international malicious actors utilizing sophisticated vectors including ransomware, targeted phishing operations, and advanced persistent threats (APTs). 
	
	The systemic impact of an unmitigated compromise extends far beyond isolated service interruptions; it introduces the risk of catastrophic capital depletion, institutional insolvency for smaller financial entities, and the complete erosion of organizational credibility. 
	
	In an unstable macro-economic landscape, a systemic failure of digital banking security risks driving the public away from formal banking institutions and toward informal, unregulated ``shadow'' financial systems, which would further cripple national economic stabilization efforts.
	
	\subsection{E-Banking Adoption Challenges and the ``Security Paradox''}
	The expansion of digital distribution channels remains significantly behind global benchmarks due to a complex mixture of psychological barriers and systemic constraints. Empirical data isolates three primary pillars retarding technology integration:
	\begin{itemize}
		\item \textbf{Heightened Risk Perception:} Widespread public apprehension regarding financial loss, driven by a fear of underlying technical system vulnerabilities and active cybercrime networks, which reinforces a cultural preference for physical cash-based transactions.
		\item \textbf{Trust and Awareness Gaps:} A profound deficit in consumer awareness regarding the operational mechanisms and tangible lifestyle advantages of electronic payments, which leaves potential users unmotivated to transition away from traditional brick-and-mortar branches.
	\end{itemize}
	
	Crucially, the empirical analysis exposes a distinct ``Security Paradox'' within the Yemeni demographic. While theoretical assumptions suggest that objective, technical data privacy vulnerabilities serve as the primary impediment to adoption, the actual data demonstrates that ``security and privacy'' metrics do not exert a statistically significant impact on a consumer’s final decision to adopt e-banking applications. Instead, adoption patterns are dominated by usability indices, general system awareness, and deep-seated cultural alignment with historical financial habits. Consequently, security compliance frameworks cannot focus exclusively on backend technical patches; they must also cultivate visible, structural trust to address this psychological barrier.
	
	\subsection{Organizational and Infrastructure Factors}
	The deployment of enterprise security controls is severely bounded by deep infrastructural deficits and institutional limitations:
	\begin{itemize}
		\item \textbf{Technical Infrastructure Constraints:} Frequent, prolonged electrical power outages and highly volatile access to standard internet connectivity present severe barriers to maintaining continuous, real-time security operations and log monitoring.
		\item \textbf{Policy Deficiencies:} Existing regulatory paradigms fail to account for the interplay between technology deployments and localized human behaviors. Local banks regularly lack formalized internal organizational policies capable of translating technical requirements into actionable, human-centric administrative workflows.
		\item \textbf{Process Immaturity:} Institutional process management is empirically verified as the single most critical factor for optimizing information security governance within Yemeni financial systems. However, ``process'' simultaneously represents the least developed component across the domestic banking architecture, highlighting a severe deficit in operational discipline.
		\item \textbf{Resource and Talent Shortages:} Financial institutions face acute shortages of specialized, certified cybersecurity personnel capable of managing complex enterprise standard requirements. This talent deficit is compounded by the high capital expenditures required for international security audits and formal certification processes, which often prove cost-prohibitive for capital-constrained local banks.
	\end{itemize}
	
	\section{PCI DSS Implementation Architecture}
	To achieve a reliable baseline of data defense, the technical architecture must secure the Cardholder Data Environment (CDE)—defined as any system component that stores, processes, or transmits cardholder account data or sensitive authentication information. The technical implementation architecture is structured around four primary deployment pillars.
	
	\subsection{Network Segmentation (VLANs and ACLs)}
	Because auditing and securing every endpoint across an entire corporate enterprise infrastructure is operationally unfeasible, the architecture leverages network segmentation as a primary scope-reduction mechanism. The broader corporate network is split into three strictly isolated zones:
	\begin{itemize}
		\item \textbf{The CDE:} The core, isolated inner network zone dedicated exclusively to the storage, processing, and transmission of cardholder account data.
		\item \textbf{The Shared Network:} Supporting systems that interface directly with the CDE but are restricted from storing sensitive cardholder data.
		\item \textbf{The Corporate LAN:} General operational environments (such as human resources, administration, and marketing departments) that are entirely blocked from accessing cardholder information.
	\end{itemize}
	
	To maintain this isolation under local operational constraints, logical segmentation is enforced at the network layer using Virtual Local Area Networks (VLANs) combined with stringent Access Control Lists (ACLs) deployed on routing and firewall perimeters. These configurations restrict traffic strictly to authorized pathways, ensuring that any lateral compromise originating within the lower-security Corporate LAN cannot easily penetrate the secure boundaries of the CDE.
	
	\subsection{Authentication and Access Control (AAA and TACACS+)}
	Physical and digital access to any system component inside or interacting with the CDE must be tightly restricted based on a documented business ``need-to-know'' parameter. The architecture enforces this through the implementation of Authentication, Authorization, and Accounting (AAA) protocols supported by Terminal Access Controller Access-Control System Plus (TACACS+). 
	
	This infrastructure ensures that any administrative connection, particularly remote configuration requests, undergoes multi-layered verification, restricts command execution to pre-authorized clearance profiles, and produces an unalterable accounting trail of all modifications made to the core security environment.
	
	\subsection{Monitoring, Logging, and Auditing}
	Continuous operational oversight is maintained by tracking and logging every access event involving network resources and cardholder files. The technical architecture mandates the centralized collection of system logs to enable rapid detection, forensic tracking, and prompt mitigation of unauthorized intrusion attempts. The monitoring apparatus must be configured to capture all user identifications, administrative actions, and perimeter modifications to satisfy the rigorous oversight criteria specified under Requirement 10 of the standard.
	
	\subsection{Data Encryption Mechanisms}
	In accordance with Requirement 4, cryptographic defenses are applied to protect account data both at rest and during transmission across open, public networks. All communications between servers, endpoints, and external databases handling Cardholder Data (CHD) must be fully encrypted using strong cryptographic protocols. This ensures that even if data packets are intercepted or captured over compromised transmission lines, the underlying financial records remain completely unreadable and useless to unauthorized parties.
	
	\section{Proposed Adaptive PCI DSS Framework for Yemeni Banks}
	The standard global PCI DSS implementation model mandates immediate, simultaneous compliance across all 12 technical and operational requirements—an approach that often fails in resource-constrained, logistically disrupted environments. To overcome these structural barriers, this section presents an adapted framework engineered specifically for the Yemeni banking ecosystem.
	
	\subsection{Framework Overview}
	The Adaptive PCI DSS Framework reframes security compliance from a rigid, ``check-the-box'' audit requirement into an incremental, risk-prioritized strategic roadmap. The model balances technical controls with localized human factors and operational constraints. It replaces the requirement for real-time high-bandwidth dependencies with localized resiliency mechanisms, optimizes human capital through targeted AI automation, and directly addresses the regional ``Security Paradox'' by emphasizing highly visible security indicators to build public trust.
	
	\subsection{Phase-Based Compliance Roadmap}
	The framework replaces simultaneous compliance with a three-tiered progression model designed to distribute capital costs and optimize limited technical resources.
	
	\subsubsection{Tier 1: Core Data Protection}
	The initial tier bypasses broader perimeter complications to focus resource allocation exclusively on securing the critical core of the CDE. Institutions must strictly deploy Requirement 3 (safeguarding stored cardholder accounts) and Requirement 4 (implementing robust cryptographic protocols). This concentration of controls ensures that even if perimeter defenses collapse or become unavailable due to localized power or network interruptions, the actual data repository remains highly secure against exploitation.
	
	\subsubsection{Tier 2: Infrastructure Resilience}
	The second tier introduces specialized structural controls to counter chronic regional electricity and internet outages. This phase requires the implementation of localized, offline storage buffers and redundant security log failover mechanisms. When primary network lines are throttled or drop entirely, security events and transaction records are safely cached in isolated, local buffers. Once connectivity is restored, these buffers automatically sync with the central security monitoring system, ensuring zero loss of audit continuity and preserving system state visibility.
	
	\subsubsection{Tier 3: Full Certification}
	The final phase builds upon the established data and infrastructure foundations to systematically implement all remaining technical and operational criteria out of the 12 global requirements. Once comprehensive security controls are running stably across the enterprise network, the bank undergoes formal international auditing and validation to secure verified PCI DSS certification.
	
	\subsection{Localized Network Segmentation and Edge Security}
	To address the regional ``Security Paradox''—where user adoption is restricted by a general perception of risk rather than technical data privacy metrics—the framework emphasizes highly visible edge security controls and robust logical boundaries. VLAN configurations are locked down using strict ``fail-safe defaults''. If a network component fails or loses state, access controls default to a completely closed posture, blocking unauthorized cross-zone traffic into the CDE. 
	
	Simultaneously, AAA and TACACS+ remote access architectures are optimized for low-bandwidth environments. By streamlining the data payloads required for authentication, authorization, and accounting, the framework ensures that administrative oversight and threat tracking do not fail or disconnect during periods of severe internet throttling.
	
	\subsection{People-First Security Pillar}
	Recognizing that organizational ``process'' is the single most influential yet least developed factor within Yemeni banks, the framework integrates a mandatory Cultural Integration Layer. This layer addresses the critical human vulnerabilities that technical controls alone cannot solve:
	\begin{itemize}
		\item \textbf{Hybrid Security Awareness Programs:} Designing internal and external training modules that pair standard PCI DSS technical concepts with counter-social engineering training. These programs are tailored to local behavioral habits to protect staff and consumers against manipulation.
		\item \textbf{Localized Policy Templates:} Generating regulatory and operational compliance documentation that translates complex, technical requirements into clear administrative workflows. These templates are customized to match the existing bureaucratic structures of Yemeni financial institutions, facilitating smooth adoption across organizational departments.
	\end{itemize}
	
	\subsection{AI-Driven Fraud Detection and Compliance Monitoring}
	To offset the acute shortage of certified cybersecurity specialists, the framework integrates advanced automation directly into the monitoring plane. By deploying specialized, artificial intelligence-driven compliance tracking tools, banks can run real-time threat detection and automated log auditing without needing large, dedicated security operations centers. 
	
	The AI layer automatically evaluates large streams of network events, flags anomalies, and isolates potential compliance infractions. This reduces the operational burden on small local technical teams while maintaining the strict oversight required by Requirement 10.
	
	\section{Comparative Discussion}
	The structural variations between the global, standard compliance baseline and the proposed adaptive framework highlight the necessity of localized security design.
	
	\subsection{Structural Comparison}
	
	\begin{table}[h!]
		\centering
		\caption{Structural Comparison between Standard and Adapted Frameworks}
		\label{tab:framework_comparison}
		\begin{tabular}{p{4cm} p{5.5cm} p{5.5cm}}
			\toprule
			\textbf{Organizational Component} & \textbf{Standard PCI DSS Baseline Framework} & \textbf{Adapted Yemeni Security Framework} \\
			\midrule
			Implementation Approach & Simultaneous deployment of all 12 requirements across the enterprise. & Tiered, phase-based progression prioritizing core data security elements first. \\
			\midrule
			Monitoring Infrastructure & Demands high-bandwidth, continuous real-time transmission architectures. & Low-bandwidth optimized design using local offline storage buffers. \\
			\midrule
			Core Risk Focus & Focuses heavily on eliminating objective technical network vulnerabilities. & Targets user risk perception and establishes visible cultural trust. \\
			\midrule
			Human Resource Model & Requires large, specialized internal security operations teams. & Leverages AI-assisted, automated monitoring tools to optimize lean teams. \\
			\bottomrule
		\end{tabular}
	\end{table}
	
	\subsection{Implementation Feasibility Under Resource Constraints}
	The standard global implementation methodology creates significant operational friction when forced onto developing financial sectors. Requiring an immediate, comprehensive security overhaul strains local banks operating under tight capital limits and severe technical deficits. 
	
	By breaking implementation into modular, phased steps, the Adaptive Framework allows local institutions to secure critical data layers immediately without exhausting their budgets. It explicitly designs around infrastructure realities—such as power drops and slow internet connections—ensuring that network drops do not break security tracking or leave data exposed. This practical approach allows Yemeni banks to steadily build towards international security standards without compromising daily operational stability.
	
	\section{Expected Benefits of the Proposed Framework}
	Adopting this localized framework provides clear strategic and technical advantages for the modernizing Yemeni financial sector:
	\begin{itemize}
		\item \textbf{Enhanced Institutional and Consumer Trust:} Operating under a structured, phased roadmap provides a verifiable ``seal of security''. This visible commitment helps lower public risk perception, bridges user trust gaps, and directly encourages wider adoption of digital banking and mobile wallets.
		\item \textbf{Systemic Strengthening of Information Security Management Systems (ISMS):} By aligning organizational workflows with the 12 core requirements of PCI DSS, banks fundamentally reinforce their entire security posture. This systemic upgrade positions local institutions to easily adopt and align with other global frameworks, such as ISO 27001.
		\item \textbf{Acceleration of Secure Digital Transformation:} Establishing a fully compliant environment creates the secure foundation necessary to launch innovative, modern financial services. It opens the door for secure internet-of-things (IoT) payments, mobile micro-finance platforms, and secure international remittance processing.
		\item \textbf{Global Financial Integration:} Aligning local security practices with international best practices helps rebuild institutional credibility on the world stage. This standard compliance is an essential step for reconnecting Yemen with international banking networks and attracting foreign investment.
	\end{itemize}
	
	\section{Chapter Summary}
	Chapter Three detailed the research methodology and structural design of the Adaptive PCI DSS Implementation Framework for the Yemeni banking sector. Grounded in a synthesis of empirical research, the chapter analyzed the unique socioeconomic conditions, infrastructure deficits, and behavioral paradoxes that define the current local banking landscape. 
	
	While standard security frameworks often struggle in developing environments due to high resource demands and rigid implementation rules, the proposed framework offers a practical path forward through a tiered, phase-based compliance model. By focusing on core data protection first, building resilience into low-bandwidth networks, introducing automated AI monitoring, and developing targeted security training, the framework effectively balances technical controls with human and organizational realities. 
	
	Ultimately, this adaptive model helps Yemeni banks move past simple, surface-level compliance checklist exercises, establishing a resilient security foundation that can drive long-term financial inclusion and national economic stability.
	
	% ==========================================
	% REFERENCES SECTION
	% ==========================================
	\begin{thebibliography}{99}
		
		\bibitem{ghallab2022}
		H. M. H. Ghallab and H. Zhu, ``Factors affecting the acceptance of electronic banking services among individuals in Yemen,'' \textit{Research Square}, 2022.
		
		\bibitem{lee2009}
		M. C. Lee, ``Factors influencing the adoption of internet banking: An integration of TAM and TPB with perceived risk and perceived benefit,'' \textit{Electronic Commerce Research and Applications}, vol. 8, no. 3, pp. 130-141, 2009.
		
		\bibitem{park1994}
		C. W. Park, D. L. Mothersbaugh, and L. Feick, ``Consumer knowledge assessment,'' \textit{Journal of Consumer Research}, vol. 21, no. 1, pp. 71-82, 1994.
		
		\bibitem{fahrizal2022}
		F. Fahrizal, A. S. Budiman, and M. R. Anuar, ``Implementasi PCI-DSS untuk Keamanan Data Kartu Pembayaran,'' \textit{Jurnal Elektro \& Informatika Swadharma (JEIS)}, vol. 2, no. 1, 2022.
		
		\bibitem{alshaibany2023}
		N. A. Al-shaibany, T. A. B. Al-sofi, and G. H. Al Gaphari, ``A Model for Enhancing the Information Security Management Systems in Yemen Banks,'' \textit{Sana'a University Journal of Applied Sciences and Technology}, vol. 1, no. 1, pp. 1-12, 2023.
		
		\bibitem{yulianto2016}
		S. Yulianto, C. Lim, and B. Soewito, ``Information security maturity model: A best practice driven approach to PCI DSS compliance,'' in \textit{2016 IEEE Region 10 Symposium (TENSYMP)}, IEEE, 2016.
		
		\bibitem{albaltah2024}
		I. A. Al-Baltah, S. Y. Al-Sultan, M. A. Al-hadi, and A. Zahary, ``Factors Influencing the Adoption of Mobile Banking Applications in Yemen Using an Extended Technology Acceptance Model,'' May 2024, doi: 10.59628/jast.v2i2.890.
		
		\bibitem{togli2022}
		М. И. Т. ўғли and K. Kato, ``Determination of the TOE Factors Influencing the Adoption of Internet Banking Services on SMEs in Yemen: A Moderated Mediation Approach,'' pp. 371-388, Jan. 2022, doi: 10.1007/978-981-19-1111-8\_30.
		
		\bibitem{factors2022}
		``Factors affecting the acceptance of electronic banking services among individuals in Yemen,'' Aug. 2022, doi: 10.21203/rs.3.rs-1863331/v1.
		
		\bibitem{alshaibany_model2023}
		N. A. Al-Shaibany, ``A Model for Enhancing the Information Security Management Systems in Yemen Banks,'' vol. 1, no. 1, Jan. 2023, doi: 10.59628/jast.vlil.14.
		
		\bibitem{alhakimi2020}
		W. Alhakimi and J. Esmail, ``The factors influencing the adoption of internet banking in Yemen,'' \textit{International Journal of Electronic Banking}, vol. 2, no. 2, p. 97, Sept. 2020, doi: 10.1504/IJEBANK.2020.10031888.
		
		\bibitem{martinez2023}
		J. G. Martinez-Navalón, M. F. Fernandez, and F. P. Alberto, ``Does privacy and ease of use influence user trust in digital banking applications in Spain and Portugal?,'' \textit{International Entrepreneurship and Management Journal}, vol. 19, no. 2, pp. 781-803, Feb. 2023, doi: 10.1007/s11365-023-00839-4.
		
		\bibitem{almuhrami2023}
		M. A. S. Al-Muhrami, N. A. Alawi, M. M. Alzubi, and A. A.-A. Al-Refaei, ``Affecting the Behavioural Intention to Use Electronic Banking Services Among Users in Yemen: Using an Extension of the Unified ...,'' 2023.
		
	\end{thebibliography}
	
\end{document}